\chapter{Methodology}
\label{ch:methodology}

\section{Overview}

\Cref{alg:main} summarises the main steps of the full pipeline. The training set is used to fit the spherical harmonic predictor and estimate the local prediction difficulty using cross-validated residuals. The separate calibration set is then used to determine the conformal quantile. Once calibration is complete, the fitted model and difficulty estimator can be applied to any query location, including held-out test points and locations on a reconstruction grid.

\begin{algorithm}[htbp]
\caption{Adaptive geodesic-neighbourhood conformal prediction for scalar fields on $S^2$}
\label{alg:main}
\begin{algorithmic}[1]
\Require Training data $\mathcal{D}_{tr}$, calibration data $\mathcal{D}_{cal}$, target level $\alpha$, degree $L$, regularisation parameters $(\lambda,k)$, neighbourhood size $k_\sigma$
\State Fit base predictor $\hat y$ on $\mathcal{D}_{tr}$ by solving \cref{eq:reg-lsq} (spherical harmonic ridge/Sobolev regression)
\State Partition $\mathcal{D}_{tr}$ into 5 folds; for each fold $r$ , refit $\hat y^{(-r)}$ on the remaining folds and record out-of-fold absolute residuals $e_i = |Y_i - \hat y^{(-r)}(X_i)|$ for all $i \in \mathcal{D}_{tr}$ in the held-out fold
\State Define $\hat\sigma(x) = \frac{1}{k_\sigma}\sum_{j \in N_{\mathrm{geo}}(x)} e_j$, where $N_{\mathrm{geo}}(x)$ is the set of $k_\sigma$ training points nearest to $x$ under the geodesic distance \cref{eq:geo-input}
\State On $\mathcal{D}_{cal}$, compute scores $s_i = |Y_i - \hat y(X_i)| / \max(\hat\sigma(X_i),\epsilon)$, and set $\hat q_\alpha$ to the split-conformal empirical quantile of $s_i$ at level $\lceil (1-\alpha)(n_{cal}+1)\rceil / n_{cal}$ 
\For{each query location $x$}
    \State Predict $\hat y(x)$ and $\hat\sigma(x)$
    \State Output interval $C(x) = [\hat y(x) - \hat q_\alpha \hat\sigma(x),\ \hat y(x) + \hat q_\alpha \hat\sigma(x)]$
\EndFor
\end{algorithmic}
\end{algorithm}

\section{Base predictor}

The underlying model is the regularised spherical harmonic regression model introduced in the background section \cref{sec:sh-background}. The model is used to reconstruct the temperature field on the sphere using a finite number of spherical harmonic basis functions. The spherical harmonic degree $L$ determines the number of basis functions included in the model and therefore indicates its complexity. The spherical harmonic coefficients are estimated using the observed temperature data to produce a point temperature prediction at any location on the sphere.
 
Regularisation is also included to control the complexity of the reconstructed field and prevent overfitting to the observed temperature data. The model uses the Sobolev-type regularisation penalty described earlier \cref{sec:sh-background}. The main model parameters, $L$-spherical harmonic degree, $\lambda$-regularisation strength, $k$-Sobolev regularisation order, determine how flexible or smooth the reconstructed temperature field is. Their values are determined through the model-selection procedure described later in \cref{sec:hyperparam-degree}.

Once $L$, $\lambda$, and $k$ are fixed, the spherical harmonic coefficients can be estimated by solving the regularised normal equation in \cref{eq:normal-eq}. Under the regularisation used, the optimisation problem is convex and the fitted coefficients have a unique solution. Therefore, in the implementation, the linear system is solved directly rather than by explicitly calculating the matrix inverse.  This gives the fitted spherical harmonic coefficients which are then used to calculate temperature predictions at different locations on the sphere.

The fitted spherical harmonic model provides the point prediction, used by the uncertainty-quantification stage, to obtain prediction residuals, which are employed to estimate the local prediction difficulty $\hat{\sigma}(x)$. The fitted predictor is then also used in conformal calibration, test prediction, and prediction at arbitrary locations. The implementation uses numerical libraries for spherical harmonic calculations, matrix operations, and solving the linear system. Further computational details are given in the software design section \cref{sec:software-design}.

\section{Difficulty estimation}
 
The difficulty estimator $\hat\sigma(x)$ represents the local prediction difficulty of the base predictor. It allows the uncertainty interval to vary according to the location. The uncertainty interval becomes large when the model tends to make more errors, while the uncertainty interval will be narrow when fewer errors are made in predictions.

The model is fitted using only the training observations, because if the model has already seen those observations, the training residuals can be artificially small and underestimate the true prediction difficulty. Instead, the proposed method should use out-of-fold residuals obtained through five-fold cross-validation on the training data, where the base predictor is fitted using four folds and the absolute prediction errors are evaluated for the remaining fold. This process is then repeated until every training observation has one out-of-fold absolute residual. These residuals are suitable because they provide a less biased estimate of prediction error, as each observation is predicted by a model that was not trained on that observation. 
 
The out-of-fold residuals can then be used to measure the local difficulty at any location. For a given location $x$, $\hat\sigma(x)$ is defined as the mean absolute residual of the $k_\sigma=20$ nearest training observations, and the neighbourhood is determined using the geodesic distance in \cref{eq:geo-input}. A small $\hat\sigma(x)$ value indicates that nearby observations have smaller errors, indicating lower local difficulty, whereas a large $\hat\sigma(x)$ means that nearby observations have larger errors, indicating higher local difficulty. Thus,  $\hat\sigma(x)$ is treated as a location-dependent measure of how difficult the reconstruction is:

\begin{equation}
\hat\sigma(x)
=
\frac{1}{k_\sigma}
\sum_{j\in N_{\mathrm{geo}}(x)} e_j,
\end{equation}

where $N_{\mathrm{geo}}(x)$ denotes the set of $k_\sigma$ training observations nearest to $x$ under the geodesic distance, and $e_j$ denotes the corresponding out-of-fold absolute residual. A small floor of $\epsilon=10^{-6}$ is applied when $\hat\sigma(x)$ is used in the normalised conformal score to avoid division by zero or extremely small values, as a numerical safeguard.

Geodesic distance is used for the neighbourhood search because latitude and longitude are angular coordinates, so directly using ordinary Euclidean distance can give misleading measures of spatial distance. This issue mostly affects higher latitudes. Geodesic distance instead measures distance along the spherical surface and provides a more appropriate measure of spatial proximity for estimating local reconstruction difficulty on $S^2$.

\section{Calibration and evaluation}
\label{sec:calibration_evaluation}

After the predictor and difficulty estimator have been fitted, the calibration set is then used to calculate the conformal threshold. The calibration follows the split conformal prediction procedure introduced earlier in \cref{sec:cp-background}.

There are two methods being compared: the standard conformal method and the adaptive conformal method. The standard conformal method uses the absolute prediction residual directly and treats the residual in the same way at every location. Therefore, it does not account for differences in prediction difficulty between locations. On the other hand, the adaptive conformal method starts by normalising the absolute prediction residual using the estimated local difficulty $\hat{\sigma}(x)$. It accounts for the fact that some locations are easier or harder to predict. The parameter $\alpha=0.1$ corresponds to a nominal coverage of 90\%. Because split conformal uses a finite-sample correction, the quantile level is adjusted using the number of calibration observations. The corrected quantile level is then used to obtain the conformal threshold from the calibration scores. With $\alpha=0.1$ corresponding to the 90\% nominal coverage level, the finite-sample-corrected quantile level is
\begin{equation}
\frac{\left\lceil(1-\alpha)(n_{\mathrm{cal}}+1)\right\rceil}
{n_{\mathrm{cal}}},
\end{equation}
The two methods are therefore evaluated as follows:

\textbf{Standard split conformal}:  
 \begin{equation}
s(x,y)=|y-\hat y(x)|
\end{equation}
This is the ordinary absolute residual. It is used to get the conformal quantile from the calibration residuals. The resulting prediction interval has the same half-width at every location, so it does not adapt to local prediction difficulty. 

\textbf{Adaptive geodesic conformal}: 
\begin{equation}
s(x,y)
=
\frac{|y-\hat{y}(x)|}
{\max\left(\hat{\sigma}(x),\epsilon\right)}.
\end{equation}
This is the proposed method described in \cref{alg:main}, where the absolute residual is normalised by dividing it by the estimated local difficulty $\hat\sigma(x)$. These two methods are evaluated on the held-out test set using four metrics:

\begin{enumerate}
\item \textbf{Marginal coverage}: this metric measures how frequent the true temperature value falls inside the prediction interval. This is evaluated against the target coverage level of 90\%. This indicates whether the method is able to achieve the target coverage across the test set.
\begin{equation}
\widehat{\mathrm{Cov}}
=
\frac{1}{n_{\mathrm{te}}}
\sum_{i\in\mathcal{D}_{\mathrm{te}}}
\mathbf{1}\{Y_i\in C(X_i)\}.
\end{equation}
where $\mathbf{1}\{Y_i\in C(X_i)\}$ denotes as an indicator that equals $1$ if the true temperature value $Y_i$ falls inside the prediction interval $C(X_i)$, and $0$ otherwise.

\item \textbf{Mean interval width}: is the average width of the prediction intervals for all observations in the test set. For the standard method, interval width is $2\hat q_\alpha$, while for the adaptive method it is $2\hat q_\alpha\hat\sigma(x)$, which varies by location. Narrower intervals are preferable, as long as the desired coverage is still maintained.

\item \textbf{Coverage by difficulty bin}: test points are grouped into equal-frequency bins based on their estimated difficulty $\hat\sigma(x)$. Three tertiles and six bins are used for comparison and more detailed analysis. The coverage is calculated separately within each bin:
\begin{equation}
    \mathrm{Cov}_b
    =
    \frac{1}{|\mathcal{B}_b|}
    \sum_{i\in\mathcal{B}_b}
    \mathbf{1}\{Y_i\in C(X_i)\},
    \qquad b=1,\ldots,B,
\end{equation}
where $\mathcal{B}_b$ represents the set of test observations in bin $b$. The standard deviation of the bin coverages is also reported to measure the consistency of coverage across bins.
\begin{equation}
    \sigma_{\mathrm{cond}}
    =
    \mathrm{std}_b(\mathrm{Cov}_b),
\end{equation}
along with the coverage of the bin that performs worst.

\item \textbf{Difficulty-estimator quality}: this is the Pearson correlation between the estimated difficulty and the absolute prediction error on the held-out test set,
\begin{equation}
    r=
    \mathrm{Corr}
    \left(
    \hat\sigma(X),
    |Y-\hat y(X)|
    \right).
\end{equation}
A higher $r$ means that the estimated difficulty better reflects the actual prediction error. This gives an indication of whether $\hat\sigma(x)$ is useful for estimating local prediction difficulty. The threshold of $r<0.15$ used by \citet{amirishahbazi2026} is also considered when judging how useful the adaptive method is here. 
\end{enumerate}

\section{Hyperparameter selection and the degree-robustness experiment}
\label{sec:hyperparam-degree}

The regularisation parameters $\lambda$ and $k$ are selected using five-fold cross-validation on the training set only. The parameter $k$ represents the Sobolev regularisation order and is also called the grad parameter used in implementation. The spherical harmonic degree is fixed at $L=26$ during the initial hyperparameter search, where this search considers $\lambda$= 21 logarithmically spaced values over the range of $10^{-8}$ to $10^{2}$ and $k\in{0,1,2,3,4}$. For every $(\lambda,k)$ combination, the model performs five-fold cross-validation and calculates the prediction error for each fold. The mean cross-validated $L^2$ prediction error is then used to select the preferred combination of regularisation parameters.
 
After the initial hyperparameter search, a separate degree-robustness experiment is then run to study whether the performance of the adaptive conformal method depends on the choice of spherical harmonic degree, $L\in{8,12,16,20,24,26,28,30}$. These different degrees represent different levels of model complexity. The purpose is to show whether the model performance remains reasonably stable as the reconstruction model becomes unstable across complexity.

For every $L$ value tested, the $\lambda$ is selected using two different criteria on the validation set. The first criterion of the $\lambda$ value is to minimise the cross-validated RMSE and pay more attention to the point reconstruction accuracy. The second criterion of the $\lambda$ value is getting the narrowest mean adaptive interval width while first maintaining a validation coverage of at least $90\%$. 

After the hyperparameters have been selected for each degree, the corresponding model is evaluated using the same calibration and test procedure described in the previous sections. \cref{sec:calibration_evaluation}. This allows comparison of the adaptive method across different spherical harmonic degrees. The two selection criteria also allow different aspects of model performance to be compared, since a parameter choice that gives a low RMSE does not necessarily produce the narrowest prediction intervals while satisfying the coverage requirement. 

\section{Software design and reproducibility}
\label{sec:software-design}

The implementation is built by three main stages: \emph{train}, \emph{reconstruct}, and \emph{degree-robustness}. Each stage has its own purpose in the overall pipeline. The \emph{train} stage fits the base predictor, selects the required hyperparameters, estimates difficulty where required, and performs conformal calibration, which produces the fitted model and calibration threshold needed for later prediction. Then, the \emph{reconstruct} stage takes the calibrated model to produce predictions on a dense spherical grid. The geographic/spatial uncertainty visualisations and diagnostics are generated in this stage described in \cref{sec:geography}. Further on, the \emph{degree-robustness} stage repeats the calibration and evaluation procedure for different spherical harmonic degrees $L$. The fitted model and calibration threshold are saved to disc using serialisation, so that computationally expensive training and calibration do not need to be repeated every time. At the same time, the reconstruction and diagnostic stages reuse the saved results without fitting the model again, which then improves both computational efficiency and reproducibility.
 
The grid-search routine \cref{sec:hyperparam-degree} is also implemented to reduce repeated computation due to some calculations that do not depend on the particular hyperparameter combination. So for each cross-validation fold, the design matrix is only constructed once and reused for all $(\lambda,k)$ combinations. The regularisation matrix $M_k$ in \cref{eq:normal-eq} is also stored for each  $k$ value and reused, this is because rebuilding it repeatedly is unnecessary when it does not depend on the observed data. In addition, the Gram matrix $A^HA$ and right-hand side $A^H\hat f$ are calculated once for each fold before looping over $\lambda$, as they can then be reused for different values of $\lambda$. The initial grid search contains 21 values of $\lambda$, 5 values of $k$, and 5 cross-validation folds, summing up to $525$ fold-level evaluations. Without caching, the design matrix could be rebuilt repeatedly which is time-consuming and computationally expensive. My implementation approach is that the design matrix only constructed once for each fold, which reduces the number of full design-matrix constructions from $525$ to $5$, with the regularised linear system being solved for each parameter combination. This reduces unnecessary computation while still evaluating all 525 parameter combinations efficiently.

The degree-robustness experiment described in \cref{sec:hyperparam-degree} also avoids repeating identical calculations. For each degree $L$, the RMSE-based and conformal-efficiency-based selection criteria may sometimes result in the same value of $\lambda$. When such a situation occurs, the calibration and test evaluation for the corresponding $(L,\lambda)$ combination is performed only once, and the results are reused for both selection criteria. For example, the two criteria produce the same selected $\lambda$ for $L=20$ and $L=26$ in \cref{tab:degree-robustness}. This avoids running the same calibration and evaluation procedure twice while retaining the results required for comparison.

% ============================================================
